% !TeX root = ../Book3.tex
\chapter*{序言}
\addcontentsline{toc}{chapter}{序言}
在一张弯曲曲面上，两点之间的距离由沿曲面行走的最短长度决定。把这个问题推广到高维流形，首先需要在每个切空间指定内积，再研究这些内积怎样随位置变化。黎曼度量由此同时产生长度、体积、联络和曲率。局部的一个二阶张量能够约束流形的直径、基本群和拓扑，也能够成为一个偏微分方程的未知量。本卷讨论这些联系，是《几何学》的第三卷。

\ProseParagraphBreak
本卷的两条线索是曲率的整体后果与几何方程的分析。沿测地线观察邻近方向怎样分开，可以把曲率条件变成距离和体积的不等式；对函数、映射、联络或子流形取变分，则得到椭圆方程和几何流。前一条线索使我们能够控制空间的形状，后一条线索使我们能够构造和改变几何结构。紧性把两者连接起来：它既用于证明极小者存在，也用于理解退化序列和奇点附近的放大极限。

\ProseParagraphBreak
本卷面向首次或第二次系统学习黎曼几何与几何分析的数学系本科生和研究生。读者需要光滑流形、切丛、微分形式、向量丛与联络的基本知识，能够使用多元微积分、有限维线性代数和常微分方程。第一卷第11—15章提供这些几何准备，讨论完整群和齐性空间时还会使用其 Lie 群及群作用内容。基本群、覆盖和 de Rham 上同调在实际应用处回顾；不必先通读第二卷的全部同伦与谱序列专题。分析部分需要测度积分、\(L^p\) 空间和 Hilbert 空间的基本训练，分布、Sobolev 截面、椭圆估计及所需的非线性函数空间工具则在本卷逐步建立。

\ProseParagraphBreak
前两编从度量进入曲率比较。第一编建立 Levi-Civita 联络、指数映射、曲率张量和子流形相容方程，并研究平行结构与黎曼分解。第二编通过能量变分、Jacobi 场和指标形式，解释曲率如何控制共轭点、距离、体积与三角形；最后进入长度空间和 Gromov--Hausdorff 收敛。这里既要辨认计算中的曲率符号，也要区分局部曲率界与完备性所带来的整体结论。

\ProseParagraphBreak
第三编集中建立后续几何分析所用的工具。Sobolev 嵌入与紧性提供弱极限，局部椭圆估计把弱解提升为正则解，全局 Fredholm 理论将方程的可解性与有限维核、余核联系起来。边值问题、Hodge 分解、Bochner 方法、热核和指标理论展示这些工具的不同用途。分析结论中的边界条件、支集、系数正则性和函数空间决定其适用范围，阅读时应把它们与所求解的方程一起保留。

\ProseParagraphBreak
第四编研究映射与子流形的能量和面积。第五编继续处理不再光滑的极限：Hausdorff 测度、可整流集、BV 函数、有限周长集、整数流和变分流形，使面积与边界在弱收敛下仍可表达。第六编把这些准备带入临界变分、Einstein 方程和几何流，研究短时存在、先验估计、长期行为及奇点模型。弱极限中的对象已经改变时，方程和正则性也需要重新说明；光滑计算、弱解存在与奇异集结构各有其具体条件。

\ProseParagraphBreak
阅读路线可以随问题选择。以黎曼几何为主，可先读前两编，再进入 Hodge、Bochner 及相应选读专题；以几何偏微分方程为主，可在第一编之后学习第11—16章，按需补齐变分和比较工具，再阅读调和映射与几何流。面积极小化的路线由第19章通向第20—23章。各编导读和章首说明给出更细的准备与依赖，星号表示可按学习目的选读，并保留该处已经注明的证明范围。

\ProseParagraphBreak
四篇附录提供不同的延伸。附录 A 检验正定性改变以后联络、测地线与因果结构的差异；附录 B 从丛联络和椭圆分析进入规范场、自旋子及四维不变量；附录 C 用输运、弱梯度和能量讨论非光滑度量测度空间；附录 D 将群的字度量、双曲边界和有限体积双曲刚性联系起来。它们分别接续正文中的具体工具，可在相应准备之后阅读。

\ProseParagraphBreak
本书用模型计算检验定义，用反例辨认条件，用完整论证说明主要方法怎样工作。习题附有解答，适合先独立尝试，再比较自己采用的条件、归一化与证明步骤。一般分类、深层紧性、规范不变量及部分奇异结构定理在相应位置准确外引；参考文献中的教材、专题专著与原始论文分别帮助回顾、深化和核查所用版本。当前正文、附录及书前书后文字均为初稿，尚未正式送审。
